\documentclass[11pt,a4paper]{article}

\usepackage{amsmath,amssymb,amsthm}
\usepackage{hyperref}
\usepackage[margin=1in]{geometry}
\usepackage{xcolor}
\usepackage{booktabs}
\usepackage{array}
\usepackage{float}
\usepackage{mathtools}

\newtheorem{theorem}{Theorem}[section]
\newtheorem{conjecture}[theorem]{Conjecture}
\newtheorem{corollary}[theorem]{Corollary}
\newtheorem{lemma}[theorem]{Lemma}
\newtheorem{proposition}[theorem]{Proposition}
\newtheorem{definition}[theorem]{Definition}
\theoremstyle{remark}
\newtheorem{remark}[theorem]{Remark}

\newcommand{\PT}{\mathbb{PT}}
\newcommand{\C}{\mathbb{C}}
\newcommand{\Z}{\mathbb{Z}}
\newcommand{\F}{\mathbb{F}}
\newcommand{\cO}{\mathcal{O}}
\newcommand{\cC}{\mathcal{C}}
\newcommand{\PP}{\mathcal{P}}
\newcommand{\Wop}[1]{W(#1)}


\title{The Cross-Context Anticommutation Theorem\\
       and the Algebraic Origin of the Kochen--Specker Obstruction\\
       {\large (Paper XVII)}}

\author{Zhou-Li Chen \\ co-nlang Research}
\date{June 2026}

\begin{document}

\maketitle

\begin{abstract}
We prove the \emph{Cross-Context Anticommutation Theorem}:
for any $K_5$ configuration of Lagrangians
$L_1, \ldots, L_5$ in $(\mathbb{F}_2^6, \omega)$
(pairwise intersecting in dimension~1),
the unique shared ray $v_{ij} = L_i \cap L_j \setminus \{0\}$
satisfies
\[
  \omega(v_{ij},\, v_{kl}) = 1
\]
for every pair $\{i,j\}, \{k,l\}$ with $\{i,j\} \cap \{k,l\} = \emptyset$
(cross-context pairs).
The proof is a pure symplectic-geometry argument in three steps:
(1)~the Fano zero-sum property (Paper~XVI)
forces all three cross-context partners of any ray $v_{ij}$
to share the same $\omega$-value;
(2)~symmetry of $\omega \bmod 2$ then forces that value to be uniform
across all 15 cross-context pairs;
(3)~a Lagrangian size argument ($10 > 2^3 - 1 = 7$) forces the value
to be~$1$, not~$0$.

As a corollary, the 15 cross-context pairs form the edge set of the
\emph{Petersen graph} $K(5,2)$, all 15 edges contributing a factor of
$-1$ to the operator product.
This gives a purely algebraic proof of the
Kochen--Specker obstruction:
$\prod_{C \in \mathbf{P}} W_C = (-1)^{15} I_8 = -I_8$
for any $K_5$ configuration of Lagrangians in $Sp(6,\mathbb{F}_2)$.
In particular, there are \emph{no even $K_5$s}:
parity $-1$ is automatic from the $K_5$ geometry,
not an additional constraint.
This resolves Open Problem~4 of Paper~XVI.
\end{abstract}


\section{Introduction}

Paper~XVI established the Weyl Product Identity $W_C = s(C) \cdot I_8$
for every context $C = \{v_1, v_2, v_3, v_4\}$,
where $s(C) = (-1)^{\beta(C)/2}$ is the context sign \cite{PaperXVI}.
The KS obstruction $\prod_{C \in \mathbf{P}} W_C = -I_8$
then followed by invoking the parity theorem of Paper~XI \cite{PaperXI},
which is itself a computational result over all 12{,}096 Mermin pentagrams.

Paper~XVI left as Open Problem~4:

\begin{quote}
\emph{Give a direct Weyl-algebraic proof that
$\prod_{C \in \mathbf{P}} W_C = -I_8$
without invoking the 12{,}096-pentagram enumeration.}
\end{quote}

This paper closes that problem.

\subsection*{Main idea}

A Mermin pentagram has 5 contexts $C_1, \ldots, C_5$ that are
\emph{pairwise adjacent}: every pair $C_i, C_j$ shares exactly one ray.
This gives the structure of a complete graph $K_5$, not a pentagon.
The 10 shared rays $v_{ij} = C_i \cap C_j$ (for $i < j$) each appear in
exactly 2 contexts; each context contains exactly 4 rays.

Expanding the product $\prod_i W_{C_i}$:
each ray $v_{ij}$ appears exactly twice (once from $C_i$, once from $C_j$).
Pairing each $W(v_{ij})$ with its duplicate:
$W(v_{ij})^2 = I_8$, but moving the second copy leftward to meet the first
picks up a factor of $(-1)^{\omega(v_{ij}, w)}$ for each operator $W(w)$
it passes.
In the standard ordering (contexts $C_1,\ldots,C_5$, rays within each
context ordered by ascending index), a cross-context pair
$\{v_{ij},v_{kl}\}$ contributes a sign of $(-1)^{\omega(v_{ij},v_{kl})}$
if and only if the two copies \emph{interleave} in the product sequence
--- i.e., the context-index intervals $[i,j]$ and $[k,l]$
overlap without nesting (Lemma~\ref{lem:pair}).
The total sign is $(-1)^{N_\mathcal{I}}$
where $N_\mathcal{I}$ is the number of anticommuting interleaving pairs.

The main theorem shows all 15 cross-context pairs anticommute.
Of these, exactly 5 interleave in the standard ordering ---
one per 4-element subset of $\{1,\ldots,5\}$ (Lemma~\ref{lem:interleave})
--- giving $(-1)^5 = -1$ directly.
Computationally verified: all 15 cross-context pairs anticommute
for all 12{,}096 Mermin pentagrams \cite{PaperXVIIsuppl}.

\subsection*{Structure of the paper}

Section~\ref{sec:k5} establishes the $K_5$ structure and the ray labeling.
Section~\ref{sec:main} states and proves the Cross-Context Anticommutation
Theorem.
Section~\ref{sec:ks} derives the KS obstruction as a corollary.
Section~\ref{sec:connections} connects the result to the Petersen graph
and to earlier papers.
Section~\ref{sec:open} states open problems.


\section{\texorpdfstring{The $K_5$ Structure of Mermin Pentagrams}{The K5 Structure of Mermin Pentagrams}}\label{sec:k5}

\begin{definition}[$K_5$ configuration]\label{def:k5}
A \emph{$K_5$ configuration} in $(V, \omega) = (\mathbb{F}_2^6, \omega)$
is a 5-tuple of Lagrangians $(L_1, \ldots, L_5)$ such that
$\dim(L_i \cap L_j) = 1$ for all $i \neq j$.
Denote by $v_{ij} = L_i \cap L_j \setminus \{0\}$
the unique non-zero shared ray.
\end{definition}

\begin{remark}[Why $K_5$, not a pentagon]
The Mermin pentagram is often drawn as a pentagon with a pentagram star
inscribed, but the adjacency structure of the 5 contexts is that of a
\emph{complete graph $K_5$}: every pair of contexts shares a ray.
The star and cycle drawings are representations, not the combinatorial
structure.
\end{remark}

The $K_5$ labeling organizes the shared rays as follows.

\begin{proposition}[$K_5$ ray statistics]\label{prop:k5rays}
A $K_5$ configuration has:
\begin{enumerate}
  \item $\binom{5}{2} = 10$ shared rays $\{v_{ij} : 1 \leq i < j \leq 5\}$.
  \item Each ray $v_{ij}$ lies in exactly 2 contexts
        ($C_i$ and $C_j$, where $C_i = \bigcup_{j \neq i} \{v_{ij}\}$).
  \item Each context $C_i$ contains exactly 4 rays
        $\{v_{ij} : j \neq i\}$.
  \item Ray pairs: $\binom{10}{2} = 45$ total,
        of which 30 are \emph{same-context pairs}
        ($\{v_{ij}, v_{ik}\}$ sharing index $i$: $5 \cdot \binom{4}{2} = 30$)
        and 15 are \emph{cross-context pairs}
        ($\{v_{ij}, v_{kl}\}$ with $\{i,j\} \cap \{k,l\} = \emptyset$:
        $\binom{10}{2} - 30 = 15$).
\end{enumerate}
\end{proposition}

\begin{proof}
Elementary counting from the $K_5$ labeling.
The 15 cross-context pairs are exactly the pairs
$\{v_{ij}, v_{kl}\}$ where $\{i,j\}$ and $\{k,l\}$ are disjoint
2-element subsets of $\{1,\ldots,5\}$; their count is
$\binom{5}{2,2,1}/2 = 15$.
\end{proof}

\begin{remark}[Same-context pairs are isotropic]\label{rem:sameiso}
For a same-context pair $\{v_{ij}, v_{ik}\}$, both rays lie in
the Lagrangian $L_i$, so $\omega(v_{ij}, v_{ik}) = 0$:
same-context pairs always commute.
The symplectic question is entirely about the 15 cross-context pairs.
\end{remark}

The product of all five context Weyl operators expands as:
\begin{equation}\label{eq:expand}
  \prod_{i=1}^5 W_{C_i}
  = \prod_{i=1}^5 \prod_{j \neq i} W(v_{ij})
  = \prod_{i < j} W(v_{ij})^2 \cdot (\text{sign from commutation moves}),
\end{equation}
where the sign arises from moving each pair $W(v_{ij})$ together
via bubble-sort transpositions.
Since $W(v_{ij})^2 = I_8$, the product equals $(-1)^{N_\mathcal{I}} \cdot I_8$
where $N_\mathcal{I}$ is the number of anticommuting cross-context pairs
that \emph{interleave} in the product ordering
(see Lemmas~\ref{lem:pair}--\ref{lem:interleave}).


\section{The Cross-Context Anticommutation Theorem}\label{sec:main}

\begin{theorem}[Cross-Context Anticommutation]\label{thm:main}
Let $(L_1, \ldots, L_5)$ be a $K_5$ configuration in
$(\mathbb{F}_2^6, \omega)$, with shared rays $v_{ij} = L_i \cap L_j \setminus \{0\}$.

For every cross-context pair $\{v_{ij}, v_{kl}\}$
(with $\{i,j\} \cap \{k,l\} = \emptyset$):
\[
  \omega(v_{ij},\, v_{kl}) = 1.
\]
\end{theorem}

\begin{proof}
We prove the theorem in three steps.

\textbf{Step~1 (Fano zero-sum forces uniform value within a row).}
Fix indices $i < j$ and pick any $m \notin \{i,j\}$.
The context $C_m$ contains the 4 rays $\{v_{mj'} : j' \neq m\}$.
By the Fano zero-sum property (Paper~XVI, Theorem~3.1):
\[
  \bigoplus_{j' \neq m} v_{mj'} = 0 \in \mathbb{F}_2^6.
\]
Applying $\omega(v_{ij}, \cdot)$ and using bilinearity:
\[
  \sum_{j' \neq m} \omega(v_{ij},\, v_{mj'}) \equiv 0 \pmod 2.
\]
Among the four terms indexed by $j' \in \{1,\ldots,5\} \setminus \{m\}$:
\begin{itemize}
  \item $j' = i$: the pair $\{v_{ij}, v_{mi}\}$ shares index $i$,
        so $v_{mi} \in L_i$ (same Lagrangian as $v_{ij}$)
        $\Rightarrow \omega(v_{ij}, v_{mi}) = 0$.
  \item $j' = j$: the pair $\{v_{ij}, v_{mj}\}$ shares index $j$,
        so $\omega(v_{ij}, v_{mj}) = 0$.
  \item $j' \notin \{i,j\}$: these are the 2 cross-context terms
        (since $\{i,j\} \cap \{m,j'\} = \emptyset$ when $j' \neq i,j,m$).
        There are exactly 2 such values of $j'$
        (namely $\{1,\ldots,5\} \setminus \{i,j,m\}$ has 2 elements).
\end{itemize}

The congruence reduces to:
\[
  \omega(v_{ij},\, v_{m,j_1'}) + \omega(v_{ij},\, v_{m,j_2'}) \equiv 0 \pmod 2,
\]
where $\{j_1', j_2'\} = \{1,\ldots,5\} \setminus \{i,j,m\}$.
Hence the two values are equal.

Applying this for all three choices $m \in \{1,\ldots,5\} \setminus \{i,j\}$:
all 3 cross-context partners of $v_{ij}$ have the \emph{same} $\omega$-value.
Call it $c_{ij} \in \{0, 1\}$.

\textbf{Step~2 (Symmetry forces a global uniform value).}
Since $\omega(v,w) = \omega(w,v) \bmod 2$ (the symplectic form is
alternating mod~2, and $-1 \equiv 1 \bmod 2$, so it is also symmetric mod~2),
$\omega(v_{ij}, v_{kl}) = \omega(v_{kl}, v_{ij})$.
Therefore $c_{ij} = c_{kl}$ for any cross-context pair $\{v_{ij}, v_{kl}\}$.

By the connectivity of the cross-context graph
(the Petersen graph $K(5,2)$, which is 3-regular and connected),
all 15 cross-context pairs share the same value $c \in \{0,1\}$.

\textbf{Step~3 (Size argument forces $c = 1$).}
Suppose for contradiction that $c = 0$.
Then all 10 rays $\{v_{ij}\}$ are mutually isotropic:
$\omega(v_{ij}, v_{kl}) = 0$ for all cross-context pairs,
and $\omega(v_{ij}, v_{ik}) = 0$ for same-context pairs (Remark~\ref{rem:sameiso}).
Hence the 10 rays span a totally isotropic subspace of $(\mathbb{F}_2^6, \omega)$.

Any totally isotropic subspace of $\mathbb{F}_2^6$ has dimension at most~3
(it is contained in a Lagrangian), hence contains at most
$2^3 - 1 = 7$ nonzero vectors.
But 10 $>$ 7 --- contradiction.

Therefore $c = 1$, and all 15 cross-context pairs satisfy
$\omega(v_{ij}, v_{kl}) = 1$. \qed
\end{proof}

\begin{remark}[Algebraic vs.\ computational]
The proof uses no case-by-case enumeration of pentagrams.
The only inputs are:
(i) the Fano zero-sum property (Paper~XVI, Theorem~3.1),
(ii) the symmetry of $\omega \bmod 2$,
(iii) the bound $\dim(\text{Lagrangian}) = 3$.
All are structural facts about $(\mathbb{F}_2^6, \omega)$.
\end{remark}

\begin{remark}[The Petersen graph]\label{rem:petersen}
The cross-context anticommutation pattern on the 10 rays is exactly
the adjacency matrix of the \emph{Petersen graph} $K(5,2)$:
\begin{itemize}
  \item Vertices: 2-element subsets $\{i,j\}$ of $\{1,\ldots,5\}$ (the 10 rays).
  \item Edges: disjoint pairs $\{\{i,j\},\{k,l\}\}$ (the 15 cross-context pairs).
\end{itemize}
Theorem~\ref{thm:main} says the Weyl anticommutation graph on the 10 rays
is \emph{exactly} the Petersen graph: edge $\Leftrightarrow$ anticommutation.
The Petersen graph is 3-regular (each ray anticommutes with exactly 3 others),
has girth~5, and is the unique $(3,5)$-cage.
\end{remark}


\section{The KS Obstruction}\label{sec:ks}

\begin{lemma}[Pairing sign]\label{lem:pair}
Let $W(u_1)\cdots W(u_{2n})$ be a product of Hermitian Pauli operators
in which each of $n$ distinct vectors $v_1,\ldots,v_n \in \mathbb{F}_2^6$
appears exactly twice.
Say that a pair $\{v_k, v_l\}$ \emph{interleaves} if the four copies
appear in the alternating pattern
$\ldots v_k \ldots v_l \ldots v_k \ldots v_l \ldots$
(or its mirror).
Then
\[
  W(u_1)\cdots W(u_{2n}) \;=\; (-1)^{|\mathcal{I}|} \cdot I,
\]
where $\mathcal{I} = \bigl\{\{v_k,v_l\} : \omega(v_k,v_l) = 1
\text{ and } \{v_k,v_l\} \text{ interleaves}\bigr\}$.
\end{lemma}

\begin{proof}
Bring the two copies of each $v_k$ together by adjacent transpositions,
processing pairs in any order.
For a non-interleaving pair $\{v_k,v_l\}$, the two copies of $v_k$
can be united passing $v_l$ an even number of times:
zero times if the pairs are separated (both copies of $v_k$ before
both copies of $v_l$, or vice versa), and two times if they are
nested (both copies of $v_k$ lie between the two copies of $v_l$).
In either case the contribution is $(-1)^0 = +1$.
For an interleaving pair, one copy of $v_k$ must pass exactly one copy
of $v_l$ to reach the other copy of $v_k$: the crossing count is odd,
contributing $(-1)^{\omega(v_k,v_l)}$.
After all pairs are united, $W(v_k)^2 = I$ cancels each pair,
leaving the accumulated sign $(-1)^{|\mathcal{I}|}$.
\end{proof}

\begin{lemma}[Interleaving count]\label{lem:interleave}
In the product $\prod_{i=1}^5 W_{C_i}$ with contexts ordered
$C_1,\ldots,C_5$ and rays within $C_i$ ordered by ascending~$j$,
exactly \emph{5} cross-context pairs interleave.
These are the pairs $\{v_{ac},\, v_{bd}\}$, one for each 4-element
subset $\{a < b < c < d\} \subset \{1,\ldots,5\}$.
\end{lemma}

\begin{proof}
The two copies of $v_{ij}$ ($i < j$) appear in block $C_i$ and
block $C_j$.
Two cross-context pairs $\{v_{ij}, v_{kl}\}$
(with $\{i,j\}\cap\{k,l\}=\emptyset$, $i<j$, $k<l$) interleave
in the sequence if and only if the context-index intervals $[i,j]$ and
$[k,l]$ overlap without nesting: $i < k < j < l$ or $k < i < l < j$.
This equivalence holds because the global positions of the two copies
of $v_{ij}$ are within blocks $C_i$ and $C_j$ respectively, so copy
ordering in the sequence is determined solely by context-block ordering.
Concretely, for indices in $\{1,\ldots,5\}$, the interleaving case
$i < k < j < l$ places the four copies in block order
$C_i, C_k, C_j, C_l$, which is the alternating pattern
(verified computationally for all 15 cross-context pairs \cite{PaperXVIIsuppl}).

For any 4-element subset $\{a < b < c < d\} \subset \{1,\ldots,5\}$,
the three cross-context pairings of its elements give intervals:
\begin{enumerate}
  \item $\{a,b\}\,|\,\{c,d\}$: $[a,b]$ and $[c,d]$ are
        \emph{separated} ($b < c$). Not interleaving.
  \item $\{a,c\}\,|\,\{b,d\}$: $a < b < c < d$ implies
        $[a,c]$ and $[b,d]$ \emph{interleave}. $\checkmark$
  \item $\{a,d\}\,|\,\{b,c\}$: $[b,c] \subset [a,d]$, \emph{nested}.
        Not interleaving.
\end{enumerate}
With $\binom{5}{4} = 5$ choices of 4-element subset, there are exactly
5 interleaving cross-context pairs.
\end{proof}

\begin{remark}[The 5 interleaving pairs form a pentagon]
\label{rem:pentagon}
The 5 interleaving pairs $\{v_{13},v_{24}\}$, $\{v_{24},v_{35}\}$,
$\{v_{35},v_{14}\}$, $\{v_{14},v_{25}\}$, $\{v_{25},v_{13}\}$
form a 5-cycle (pentagon) in the Petersen graph $K(5,2)$
(see Remark~\ref{rem:petersen}).
Since the 5-cycle has an odd number of edges and all 5 pairs
satisfy $\omega = 1$, the sign is $(-1)^5 = -1$.
\end{remark}

\begin{corollary}[Algebraic KS Obstruction]\label{cor:ks}
For any $K_5$ configuration of Lagrangians in $(\mathbb{F}_2^6, \omega)$,
\[
  \prod_{i=1}^5 W_{C_i} \;=\; -I_8.
\]
In particular, $\prod_{C \in \mathbf{P}} s(C) = -1$ for every Mermin pentagram.
\end{corollary}

\begin{proof}
By Lemma~\ref{lem:interleave}, in the standard product ordering,
exactly 5 cross-context pairs interleave.
By Theorem~\ref{thm:main}, all 15 cross-context pairs (and hence
all 5 interleaving ones) satisfy $\omega(v_{ij},v_{kl}) = 1$.
By Lemma~\ref{lem:pair}:
\[
  \prod_{i=1}^5 W_{C_i}
  \;=\; (-1)^5 \cdot I_8
  \;=\; -I_8.
\]
Combined with $W_{C_i} = s(C_i) \cdot I_8$ (Paper~XVI, Theorem~4.1):
$\prod_i s(C_i) = -1$. \qed
\end{proof}

\begin{remark}[Comparison with Paper~XVI]
Paper~XVI derived $\prod_C W_C = -I_8$ by:
(a) proving $W_C = s(C) I_8$ algebraically (Theorem~4.1), and
(b) citing the parity theorem of Paper~XI ($\prod s(C) = -1$)
    which relied on a 12{,}096-pentagram enumeration.

Corollary~\ref{cor:ks} replaces step~(b) with the algebraic
Theorem~\ref{thm:main}.
The result is now a pure theorem of symplectic geometry over $\mathbb{F}_2$,
with no computer enumeration in the logical chain.
\end{remark}

\begin{corollary}[No even $K_5$s]\label{cor:noeven}
There are no $K_5$ configurations of Lagrangians in $Sp(6, \mathbb{F}_2)$
with $\prod_{i=1}^5 s(C_i) = +1$.
The parity $\prod s(C_i) = -1$ is \emph{automatic} from the $K_5$ geometry.
Computationally: all 12{,}096 $K_5$ configurations have parity $-1$
\cite{PaperXVIIsuppl}.
\end{corollary}

\begin{proof}
Immediate from Corollary~\ref{cor:ks}: every $K_5$ configuration
satisfies $\prod s(C_i) = -1$, so no $K_5$ can have product $+1$.
\end{proof}

\begin{remark}[Consequence for the KS theorem]
Corollary~\ref{cor:noeven} shows that the Mermin pentagram condition
``is a pentagram'' (which requires $\prod s(C) = -1$ for the
no-coloring argument) is automatically satisfied by \emph{every}
$K_5$ configuration of Lagrangians in $Sp(6,\mathbb{F}_2)$.
The KS obstruction is not a property of a special subset of $K_5$s;
it is a universal property of the $K_5$ geometry.
\end{remark}


\section{Connections to Previous Papers}\label{sec:connections}

\subsection{Resolution of Paper~XVI Open Problem~4}

Theorem~\ref{thm:main} and Corollary~\ref{cor:ks} together close
Open Problem~4 of Paper~XVI:

\begin{quote}
\emph{Give a direct Weyl-algebraic proof that
$\prod_{C \in \mathbf{P}} W_C = -I_8$
without invoking the 12{,}096-pentagram enumeration.}
\end{quote}

The algebraic chain is now complete:
\[
  \underbrace{\text{Fano zero-sum}}_{\text{Paper~XVI, Thm~3.1}}
  \;\Rightarrow\;
  \underbrace{N_{\mathrm{anti}} = 15}_{\text{Thm~\ref{thm:main}}}
  \;\Rightarrow\;
  \underbrace{N_\mathcal{I} = 5 \text{ (odd)}}_{\text{Lem~\ref{lem:interleave}}}
  \;\Rightarrow\;
  \underbrace{\prod W_C = (-1)^5 I_8 = -I_8}_{\text{Cor~\ref{cor:ks}}}.
\]

\subsection{The T-Vector Theorem (Paper~XII)}

Paper~XII established the existence of a \emph{T-vector}
$T \in \mathbb{F}_2^6$ satisfying $\omega(T, v) = 1$ for all 10 rays
of any Mermin pentagram \cite{PaperXII}.

Theorem~\ref{thm:main} provides context: since all 15 cross-context
pairs anticommute, the 10 rays are not mutually isotropic, confirming
that the system $\omega(T,v)=1$ for all 10 rays is non-trivial.
The spanning result --- that the 10 rays span all of $\mathbb{F}_2^6$
--- was established in Paper~X (Theorem~6.4: the 10 shared operators
form a 10-point cap in $PG(5,\mathbb{F}_2)$, which spans $\mathbb{F}_2^6$)
\cite{PaperX}.
Computationally, $T$ is the unique solution to a rank-6 linear system
$\{AT = \mathbf{1}\}$ over $\mathbb{F}_2$, and uniqueness holds for all
12{,}096 pentagrams \cite{PaperXVIIsuppl}.

\subsection{The Parity Theorem (Paper~XI)}

Paper~XI established $\prod s(C) = -1$ computationally
(all 12{,}096 pentagrams) \cite{PaperXI}.
Corollary~\ref{cor:ks} now provides a purely algebraic proof.
The computational result is subsumed.

\subsection{Updated Resolution Table}

\begin{table}[H]
\centering
\caption{Resolution of open problems from Papers~XI--XVI.}
\label{tab:resolution}
\begin{tabular}{p{1.4cm}p{5.0cm}p{5.4cm}}
\toprule
Paper & Open Problem & Resolution \\
\midrule
XI & Algebraic proof of $\prod s(C) = -1$
   & \textbf{Paper~XVII:} Cross-Context Anticommutation Theorem \\
\addlinespace
XII & T-vector existence
   & \textbf{Paper~XVII:} Structural (10 rays span $\mathbb{F}_2^6$,
     rank-6 system) \\
\addlinespace
XIII--XV & Maslov / Weil / cocycle attempts
   & Resolved by Paper~XVI (Weyl Product Identity) \\
\addlinespace
XVI OP~4 & Algebraic $\prod W_C = -I_8$
   & \textbf{Paper~XVII:} Cor~\ref{cor:ks} (this paper) \\
\bottomrule
\end{tabular}
\end{table}


\section{Open Problems}\label{sec:open}

\begin{enumerate}

\item \textbf{$n$-qubit generalization.}
      Does Theorem~\ref{thm:main} hold in $Sp(2n, \mathbb{F}_2)$
      for arbitrary $n \geq 2$?
      The Fano zero-sum and bilinearity arguments (Steps~1--2) generalize
      directly; only Step~3 (the size argument) depends on the specific
      bound $10 > 2^3 - 1 = 7$.
      For $n \geq 4$: a Lagrangian in $\mathbb{F}_2^{2n}$ has $2^n - 1$
      nonzero vectors, and $2^4 - 1 = 15 \geq 10$, so the size argument
      \emph{fails}.
      Either the theorem requires a different proof for $n \geq 4$,
      or there exist $K_5$ configurations with $c = 0$.

\item \textbf{Petersen graph and KS no-coloring.}
      The Petersen graph is famously non-3-colorable (it has chromatic
      number~4), and the KS theorem asserts that the 10-ray system
      has no binary valuation (coloring) consistent with the measurement
      constraints.
      Is there a direct combinatorial proof connecting the
      non-3-colorability of the Petersen graph to the KS no-coloring
      result for the 10 rays?

\item \textbf{Other symplectic geometries.}
      Does an analog of Theorem~\ref{thm:main} hold in $Sp(6, \mathbb{F}_q)$
      for odd primes $q$, or in other finite classical geometries?
      The Fano zero-sum property is specific to $\mathbb{F}_2$
      (it uses the fact that $|L| = 2^3$ forces $\bigoplus_{v \in L} v = 0$);
      its analogue over $\mathbb{F}_q$ for $q > 2$ may require different
      combinatorial input.

\item \textbf{Weil representation interpretation.}
      The context-level sign $s(C) = (-1)^{\beta(C)/2}$ is realised
      as the Pauli product $W_C$ (Paper~XVI).
      The pentagram-level product $\prod s(C) = -1$
      is realised as the cross-context anticommutation count (this paper).
      Is there a single Weil-representation-theoretic statement
      that encodes both levels simultaneously?

\end{enumerate}


\appendix
\setcounter{table}{4}
\renewcommand{\thetable}{A\arabic{table}}
\renewcommand{\theHtable}{A\arabic{table}}
\section{Rigor Self-Assessment}

\begin{table}[H]
\centering
\caption{Rigor classification for all main results.}
\label{tab:rigor}
\begin{tabular}{p{6.2cm}lp{4.0cm}}
\toprule
Result & Type & Scope \\
\midrule
Prop~\ref{prop:k5rays} ($K_5$ ray statistics)
  & Algebraic & Combinatorial counting \\
Thm~\ref{thm:main} (Cross-Context Anticommutation)
  & Algebraic & 3-step pure proof \\
Cor~\ref{cor:ks} ($\prod W_C = -I_8$)
  & Algebraic + Cite & Thm~\ref{thm:main} + Paper~XVI Thm~4.1 \\
Cor~\ref{cor:noeven} (No even $K_5$s)
  & Algebraic + Comp & Follows from Cor~\ref{cor:ks};
                       12{,}096/12{,}096 verified \cite{PaperXVIIsuppl} \\
Rem~\ref{rem:petersen} (Petersen graph)
  & Algebraic & Combinatorial identity \\
T-vector uniqueness (\S\ref{sec:connections})
  & Computational & 12{,}096/12{,}096 verified \cite{PaperXVIIsuppl} \\
\bottomrule
\end{tabular}
\end{table}


\begin{thebibliography}{99}

\bibitem[PaperX]{PaperX}
Z.-L. Chen,
``The Equiangular Characterization of Mermin Pentagrams:
Symplectic Embedding, 10-Ray Cap, and the Failure of the
Collinearity--Obstruction Dictionary,''
\textit{Zenodo, DOI: 10.5281/zenodo.20476659} (2026).

\bibitem[PaperXI]{PaperXI}
Z.-L. Chen,
``Quadratic Refinement and the Parity of Mermin Pentagrams:
The $\beta$-Formula, Geometric Decomposition, and the
Kochen--Specker Obstruction from Equiangular Geometry,''
\textit{Zenodo, DOI: 10.5281/zenodo.20482595} (2026).

\bibitem[PaperXII]{PaperXII}
Z.-L. Chen,
``The T-Vector Theorem: Symplectic Characteristic Elements
and the Kochen--Specker Obstruction,''
\textit{Zenodo, DOI: 10.5281/zenodo.20490118} (2026).

\bibitem[PaperXIII]{PaperXIII}
Z.-L. Chen,
``The Maslov Index and the Kochen--Specker Obstruction:
Negative Results, the $k$-Profile Theorem, and Orbit~4 Anomaly,''
\textit{Zenodo, DOI: 10.5281/zenodo.20496513} (2026).

\bibitem[PaperXIV]{PaperXIV}
Z.-L. Chen,
``Stabilizer Algebra and the $k$-Profile Theorem for Mermin Pentagrams,''
\textit{Zenodo, DOI: 10.5281/zenodo.20502868} (2026).

\bibitem[PaperXV]{PaperXV}
Z.-L. Chen,
``The Weil Representation and the $S_3$ Lifting Classification
for Mermin Pentagrams,''
\textit{Zenodo, DOI: 10.5281/zenodo.20519654} (2026).

\bibitem[PaperXVI]{PaperXVI}
Z.-L. Chen,
``The Weyl Product Identity and the Algebraic Origin
of the Kochen--Specker Obstruction,''
\textit{Zenodo, DOI: 10.5281/zenodo.20519733} (2026).

\bibitem[PaperXVIIsuppl]{PaperXVIIsuppl}
Z.-L. Chen,
``Supplementary computational scripts for Paper~XVII,''
\textit{Zenodo, DOI: 10.5281/zenodo.20528739} (2026).

\bibitem[Mermin]{Mermin}
N.D. Mermin,
``Hidden variables and the two theorems of John Bell,''
\textit{Rev. Mod. Phys.} \textbf{65} (1993), 803--815.

\bibitem[Petersen]{Petersen}
J. Petersen,
``Sur le th\'eor\`eme de Tait,''
\textit{L'Interm\'ediaire des Math\'ematiciens} \textbf{5} (1898), 225--227.

\end{thebibliography}

\end{document}
